\begin{pastexam}[label={px:ccu-109-swa-02}]{中正資工所甲組 109 學年度 第 5 題}
（程式）

(10\%) Please write a function to perform linked list insertion. Assume that each node has a name
field with at most 32 bytes, and a link field. You have to declare the data structure for the node.
The inserted node will become the first node of the linked list. The return value of the function
is a pointer to the head of the linked list.
\end{pastexam}

\begin{pastsoln}
「name field with at most 32 bytes」採用的讀法：\texttt{name} 欄位本身佔 $32$ 個 byte，C 字串要留一格
給結尾的 \verb|'\0'|，所以名字最多 $31$ 個字元。 如果要讓名字本身最多 $32$ 個字元，把陣列宣告成
\texttt{name[33]} 即可，其他不變。

做法就是資結書虛擬碼~5.2 裡「接在最前面」的情形（資結書第~5.1~節）：新節點的 \texttt{link} 指向原來
的第一個節點，再把新節點當成新的串列頭回傳。 函式拿舊的 \texttt{head} 當參數、回傳新的
\texttt{head}，呼叫端寫 \texttt{head = insert\_front(head, "Amy");}。

\begin{minted}[fontsize=\small, style=perldoc]{c}
#include <stdlib.h>
#include <string.h>

struct node {
    char name[32];          /* 名字，含結尾的 '\0' 共 32 bytes */
    struct node *link;      /* 下一個節點；最後一個是 NULL */
};

/* 把 name 插在串列最前面，回傳新的串列頭 */
struct node *insert_front(struct node *head, const char *name)
{
    struct node *p = malloc(sizeof(struct node));
    if (p == NULL)          /* 配置失敗：串列不變 */
        return head;
    strncpy(p->name, name, sizeof(p->name) - 1);
    p->name[sizeof(p->name) - 1] = '\0';   /* 名字太長時截斷，仍保證有結尾 */
    p->link = head;         /* 新節點接在原來的第一個節點前面 */
    return p;               /* 新節點成為串列頭 */
}
\end{minted}

\textbf{正確性}。 設插入前串列是 $x_1 \to x_2 \to \cdots \to x_k \to \texttt{NULL}$，\texttt{head} 指向
$x_1$。 執行 \texttt{p->link = head} 之後，從 \texttt{p} 出發依序是 $p, x_1, \ldots, x_k$，原來的節點
一個都沒少、順序不變；回傳 \texttt{p}，新節點就是第一個。 空串列（\texttt{head == NULL}）也一樣：
\texttt{p->link} 是 \texttt{NULL}，串列只有新節點一個，不用另外處理。 兩行的順序不能倒過來：先回傳
或先讓 \texttt{head} 改指 \texttt{p}，就找不到原來的 $x_1$ 了。

\textbf{複製字串}。 \texttt{strncpy} 最多複製 $31$ 個字元；來源比較長時它不會自己補
\verb|'\0'|，所以最後一格手動設成 \verb|'\0'|。 直接寫 \texttt{strcpy(p->name, name)}，名字超過
$31$ 個字元就會寫出陣列範圍。

\textbf{複雜度}。 配置一個節點、複製最多 $32$ 個 byte、改一個指標，和串列長度無關，$O(1)$。

例如從空串列依序插入 \texttt{"apple"}、\texttt{"banana"}，串列變成
\texttt{banana} $\to$ \texttt{apple} $\to$ \texttt{NULL}。
\end{pastsoln}
